\section{Matter: the Kohn-Sham analysis}
\label{sec:ks}

This section determines the states of uniform matter made of the fermions of the model when the
self-interaction acts between them. The Vorton-Scalar is set aside here ($\yVS=0$, so that the mass
$m$ of the fermions in vacuum is the bare Dirac mass); Sec.~\ref{sec:cjtcontact} restores it.

\subsection{The method}
\label{sec:ksmethod}

In the method of Kohn and Sham \cite{KohnSham1965} a system of interacting fermions is replaced by
independent fermions that move in fields which their own densities determine. The energy is written
as a function of the densities; the fields are its derivatives; and the densities are computed from
the occupied states in those fields, until both agree. \SIG{The method is set out here for matter of
one kind of particle, whose densities are two numbers}{For uniform matter of the present model the
densities are two numbers (Sec.~\ref{sec:hf})}: the number density $n=\langle\psi^\dagger\psi\rangle$
and the scalar density $\nS=\langle\bar\psi\psi\rangle$. \SIG{(Section~\ref{sec:ksreduced} treats the
two halves of the field together.) }{}The interaction energy is
$\tfrac12\GV n^{2}+\tfrac12\GS\nS^{2}$ with a scalar coupling $\GS$ and a vector coupling $\GV$, and
each particle obeys the Dirac equation with the Hamiltonian matrix
\begin{equation}
\begin{aligned}
\hh_{\mathrm{KS}}&=\sum_j\alpha^{j}k_j+\beta\,\mst+V_0,\\
\mst&=m+\GS\,\nS,\qquad V_0=\GV\,n .
\end{aligned}
\label{eq:hKS}
\end{equation}
The scalar field changes the mass of the particles in matter from $m$ to $\mst$; the vector field
shifts all energies by $V_0$. A negative $\GS$ is an attraction (it lowers the mass and the energy)
and a positive $\GV$ a repulsion; this is the structure of the mean-field theory of dense matter
with a scalar and a vector field \cite{Walecka1974}. With the couplings of Eq.~(\ref{eq:GHF}) the exchange terms are
contained exactly, because for a contact interaction they are a function of the same two densities.

A particle of the momentum $\vec k$ has the energy $E^{*}(k)+V_0$ with
$E^{*}(k)=(k^{2}+\mstq)^{1/2}$. The occupied states are those with $|\vec k|<\kF$, \SIG{$\nu$ for each
momentum ($\nu=8$ if all eight internal states of a particle are filled, $\nu=4$ if the particles of
one half of the field are present alone)}{eight for each
momentum ($\nu=8$)}, and averages over them are written
$\langle f\rangle=\nu\int d^{D}k\,f/(2\pi)^{D}$, where $D$ is the number of space dimensions in which
the particles have momentum. Then
\begin{equation}
n=\langle1\rangle=\frac{\nu\,V_D\,\kF^{D}}{(2\pi)^{D}},\qquad\nS=\Big\langle\frac{\mst}{E^{*}}\Big\rangle,
\label{eq:densities}
\end{equation}
with the volume $V_D=\pi^{D/2}/\Gamma_{\mathrm E}(D/2+1)$ of the unit ball ($\Gamma_{\mathrm E}$ is
Euler's gamma function). The second of Eqs.~(\ref{eq:hKS}) with the second of
Eqs.~(\ref{eq:densities}) is the \emph{gap equation} for $\mst$. The energy density is
\begin{equation}
\varepsilon=\langle E^{*}\rangle+\tfrac12\GV\,n^{2}-\frac{(\mst-m)^{2}}{2\GS},
\label{eq:epsKS}
\end{equation}
which is the kinetic and mass energy of the particles plus the interaction energy [the first term
counts the mass as $\mst$ instead of $m$, which the last term corrects]; it is stationary with
respect to $\mst$ exactly when the gap equation holds. The pressure along the coordinate $xj$ is
\begin{equation}
p_j=\Big\langle\frac{k_j^{2}}{E^{*}}\Big\rangle+\pint,\qquad\pint=\tfrac12\GV\,n^{2}+\frac{(\mst-m)^{2}}{2\GS}.
\label{eq:pKS}
\end{equation}
The first term is the momentum that the particles carry across a surface; the \emph{interaction
pressure} $\pint$ is the same in all directions. \emph{Along a coordinate on which the fields do not
depend the particles have no momentum, and the pressure is $\pint$ alone}; this remark is used for
the deflating coordinates below. The chemical potential is $\mu=E^{*}(\kF)+V_0$, and for an
isotropic state $p=\mu n-\varepsilon$. Matter is \emph{self-bound} if its energy per particle,
$\varepsilon/n$, is below the mass $m$ at some density; the density at which $\varepsilon/n$ is
smallest is the \emph{saturation} density, where the pressure vanishes and
$\mu=\varepsilon/n$ \cite{HugenholtzVanHove1958}.

Two limitations are stated here once. The filled states of negative energy are not counted (the
\emph{no-sea approximation}); Sec.~\ref{sec:cjtcontact} counts them. And the state is the one of
lowest energy among the solutions of the gap equation: without the filled negative-energy states the
energy as a function of an unrestricted trial mass has no lowest value, because a large negative
trial mass places particles in states of negative rest energy.

%%<1671>
\subsection{Flat spacetime: the reference calculation}
\label{sec:ksflat}

On the flat $7+1$ spacetime the particles have momentum along all seven space coordinates, $D=7$,
and a coupling has the dimension of a length to the sixth power. The reference example has
$\GS=-10^{5}$ and $\GV=+10^{4}$ in units $m=1$. Table~\ref{tab:ksflat} gives its saturation point,
computed by the program of this paper, together with the values of a first calculation with a
notebook in the Wolfram Language; an independent script reproduces the same numbers a third time.
Figure~\ref{fig:ksflat} shows the equation of state.

\begin{table*}
\caption{\label{tab:ksflat}Uniform matter on the flat $7+1$ spacetime at saturation, for
$\GS=-10^{5}$, $\GV=+10^{4}$ (units $m=1$, seven space dimensions): the program of this paper and
the first calculation with a notebook.}
\begin{ruledtabular}
\begin{tabular}{lll}
quantity & this program & notebook \\
\hline
{{ROWS:tbl_ks_flat}}
\end{tabular}
\end{ruledtabular}
\end{table*}

\begin{figure*}
\includegraphics[width=\textwidth]{fig_ks_flat}
\caption{\label{fig:ksflat}Uniform matter on the flat $7+1$ spacetime (seven space dimensions; units
$m=1$) for two pairs of mean-field couplings. (a) Energy per particle minus $m$; the circle marks the
saturation point of Table~\ref{tab:ksflat}. (b) Mass of the particles in matter. (c) Pressure
divided by the energy density.}
\end{figure*}

The matter is self-bound: at $\kF={{GENM:ks_kf_sat}}\,m$ the energy per particle is
${{GENM:ks_ebind}}\,m$ below the mass, and the mass of the particles in it is
$\mst={{GENM:ks_mstar}}\,m$. The identity $p=\mu n-\varepsilon$ holds along the whole equation of state
to $10^{-15}$ of $\varepsilon$, and at saturation $\mu=\varepsilon/n$.

One remark concerns the relation between these two numbers and the couplings of the Lagrangian. If
the Lagrangian has $g_0=-10^{5}$, $g_1=+10^{4}$ and no other coupling, Eq.~(\ref{eq:GHF}) gives
$\GS^{\mathrm{HF}}=-98750$ and $\GV^{\mathrm{HF}}=+20000$: the exchange terms double the vector
repulsion. For that pair the energy per particle is above $m$ at all densities [dashed curves of
Fig.~\ref{fig:ksflat}; its local minimum lies ${{GENM:ks_hf_localmin}}\,m$ above $m$, at
$\kF={{GENM:ks_hf_localmin_kf}}\,m$]: \emph{uniform matter is not self-bound.} The reference example is therefore to be read as $\GS=-10^{5}$, $\GV=+10^{4}$ for the
mean-field couplings themselves, which corresponds to $g_0=-105341.6$, $g_1=+2484.5$ if no other
coupling of the Lagrangian differs from zero.

\subsection{Matter with momentum along the deflating coordinates}
\label{sec:ks71}

Now let $x5,x6,x7$ deflate, and consider matter whose particles have momentum along all seven space
coordinates. The linear Dirac equation conserves the wavenumbers $k_t$ of a particle
(Sec.~\ref{sec:dirac}); its proper momentum along $xt$ is $k_t/b_t$ and grows as the coordinate
shrinks. If the particles do not collide, the occupied states remain the ball $|\vec k|<\kF$ of
wavenumbers that they filled at a reference time, and the occupied proper momenta form an
\emph{ellipsoid} that is stretched along $x5,x6,x7$ by the factors
$b_t^{\mathrm{ref}}/b_t=e^{\rt(x4-x4_{\mathrm{ref}})}$. The number density grows as the volume
shrinks,
\begin{equation}
n=n_{\mathrm{ref}}\,e^{\td},
\end{equation}
and the equations of Sec.~\ref{sec:ksmethod} hold with the averages taken over the ellipsoid; the
pressure along a deflating coordinate is $p_t=\langle(k_t/b_t)^{2}/E^{*}\rangle+\pint$. The energy
density obeys
\begin{equation}
\frac{d\varepsilon}{dx4}=\sum_{t}\rt\,(\varepsilon+p_t),
\label{eq:conservation}
\end{equation}
the statement that the covariant divergence of the energy-momentum tensor vanishes; the program
confirms it to $6\times10^{-9}$ and the independent script to $2\times10^{-8}$.

Figure~\ref{fig:ksdefl} shows the history of the saturated matter of Table~\ref{tab:ksflat}, taken as
the state at the reference time, for three equal rates and for rates in the ratio $1:2:3$.

\begin{figure*}
\includegraphics[width=\textwidth]{fig_ks_deflating}
\caption{\label{fig:ksdefl}Matter of the $7+1$ model with momentum along all seven space coordinates
while $x5,x6,x7$ deflate, without collisions. At the clock reading $\td=0$ it is the saturated matter
of Table~\ref{tab:ksflat}. (a) Mass of the particles in matter. (b) Energy per particle minus $m$:
for the collisionless state, and for the state of lowest energy at the same density (``relaxed'').
(c) Pressures along the flat coordinates and along the three deflating ones, divided by the energy
density, for rates in the ratio $1:2:3$.}
\end{figure*}

\emph{The mass of the particles falls.} It reaches half its value at the reference time at
$\td={{GENM:ks71_half_eq}}$ (equal rates) or ${{GENM:ks71_half_123}}$ (rates $1:2:3$).

\emph{Matter is self-bound only for a limited time.} The energy per particle is below $m$ only for
${{GENM:ks71_lo_eq}}<\td<{{GENM:ks71_hi_eq}}$ (equal rates;
${{GENM:ks71_lo_123}}<\td<{{GENM:ks71_hi_123}}$ for $1:2:3$): about one e-folding of the deflating
volume. Before that the matter is too dilute to bind, afterwards the vector repulsion
$\tfrac12\GV n^{2}$ dominates; at $\td=4$ it is {{GEN:ks71_vec_eq}} per cent of the energy density,
which then grows as $e^{2\td}$.

\emph{The pressure becomes anisotropic}, larger along the coordinates that deflate faster
[Fig.~\ref{fig:ksdefl}(c)].

\emph{Collisions would release heat.} The collisionless state has more energy than the state of
lowest energy at the same density, a ball of occupied states; the difference reaches
${{GENM:ks71_heat}}\,m$ per particle at $\td=4$. Whether collisions are frequent enough to release it
is not examined here.
%%</1671>

\subsection{Fields that do not depend on the deflating coordinates}
\label{sec:ksreduced}

%%<1644>
On the $4+4$ spacetime matter exists only in the quantized sector: a mode that depends on a
deflating time is extinguished (Sec.~\ref{sec:modes}). The particles then have momentum along
$x0,\dots,x3$, so that $D=4$, or, if the fields do not depend on the hidden coordinate $x0$ either,
along ordinary space, $D=3$. By Eq.~(\ref{eq:couplings}) the equations of flat spacetime hold for the
weighted field at each moment with the couplings of that moment: every coupling $g_r$, and with it
every mean-field coupling of Sec.~\ref{sec:hf}, grows as $e^{\td}$, while the mass $m$ and the
density of particles per unit of the volume of $x0,\dots,x3$ do not change. The ratios of the
couplings are fixed. The deflation therefore carries the matter through a one-parameter sequence of
interactions of growing strength.

\emph{The equations of two halves.} The field has two halves (Sec.~\ref{sec:forms}), and the
interaction energy (\ref{eq:EHF}) depends on the densities of both. The particles of the half
$\Wg=\pm1$ have the mass $\mstpm$ and the energy shift $V_{0\pm}$,
\begin{equation}
\begin{aligned}
\mstpm&=m+\GStot\,\nS\pm\GSdif\,n_{SW},\\
V_{0\pm}&=\GVtot\,n\pm\GVdif\,n_W ,
\end{aligned}
\label{eq:hKShalves}
\end{equation}
and fill the states up to the Fermi momenta $k_{F\pm}$, four for each momentum. With
$E^{*}_\pm(k)=(k^{2}+m_\pm^{*2})^{1/2}$,
\begin{equation}
\begin{aligned}
n_\pm&=4\!\int_{|\vec k|<k_{F\pm}}\!\frac{d^{D}k}{(2\pi)^{D}},\\
n_{S\pm}&=4\!\int_{|\vec k|<k_{F\pm}}\!\frac{d^{D}k}{(2\pi)^{D}}\,\frac{\mstpm}{E^{*}_\pm},
\end{aligned}
\label{eq:denshalves}
\end{equation}
\begin{equation}
\begin{aligned}
\varepsilon={}&\sum_\pm4\!\int_{|\vec k|<k_{F\pm}}\!\frac{d^{D}k}{(2\pi)^{D}}\,E^{*}_\pm
+\tfrac12\big(\GVtot\,n^{2}+\GVdif\,n_W^{2}\big)\\
&-\tfrac12\big(\GStot\,\nS^{2}+\GSdif\,n_{SW}^{2}\big),
\end{aligned}
\label{eq:epshalves}
\end{equation}
and the pressure along $x0,\dots,x3$ is the sum of the kinetic terms of Eq.~(\ref{eq:pKS}) for the two
halves and of the interaction pressure
$\pint=\tfrac12(\GVtot n^{2}+\GVdif n_W^{2})+\tfrac12(\GStot\nS^{2}+\GSdif n_{SW}^{2})$, which is the
whole pressure along the deflating coordinates. The first of Eqs.~(\ref{eq:hKShalves}) with the second
of Eqs.~(\ref{eq:denshalves}) are two coupled gap equations for $m^{*}_+$ and $m^{*}_-$. If one half is
empty, or if both halves are equally filled, these equations are those of Sec.~\ref{sec:ksmethod} for
one kind of particle: with $\nu=4$ and the couplings $\GSsame$, $\GVsame$ for one filled half, and
with $\nu=8$ and the couplings $\GStot$, $\GVtot$ for equal halves. Because the interaction can move
two particles from one half to the other (Sec.~\ref{sec:forms}), the numbers of the two halves are not
fixed separately: at a given total density $n$ the state is the one of lowest energy with respect to
the \emph{polarization}
\begin{equation}
\mathrm{pol}_W=\frac{n_+-n_-}{n},
\label{eq:pol}
\end{equation}
between 0 (equal halves) and 1 (one half filled; the sign is free by the symmetry between the
halves). Among the solutions of the gap equations those with non-negative masses are taken, which
continue the dilute gas. Such a solution is a regular function of the density and of the couplings
while the eigenvalues of the matrix $\mathsf R_{\mathrm{gap}}$ of the derivatives of the right-hand
sides of the two gap equations with respect to the two masses stay below 1; they are
$2\,\GStot\,\partial n_{S+}/\partial m^{*}_+$ and $2\,\GSdif\,\partial n_{S+}/\partial m^{*}_+$ for equal
halves, and $\GSsame\,\partial n_{S+}/\partial m^{*}_+$ and 0 for one filled half.

The numbers below are given for two sets of couplings of the Lagrangian, each with one strength
$g>0$ that grows as $e^{\td}$: \emph{set A}, a scalar attraction alone, $g_0=-g$; and \emph{set B},
$g_0=-g$ with a vector repulsion one tenth as large, $g_1=+0.1\,g$; all other couplings zero. For
these Table~\ref{tab:kscoup} gives the mean-field couplings of Eqs.~(\ref{eq:EHF}) and
(\ref{eq:Esame}).

\begin{table*}
\caption{\label{tab:kscoup}The mean-field couplings of the two halves of the field for the sets A
and B of couplings of the Lagrangian, in the two forms, in units of the strength $g$ (exact).}
\begin{ruledtabular}
\begin{tabular}{cccccccccc}
set & form & $\GVtot$ & $\GStot$ & $\GVdif$ & $\GSdif$ & $\GVsame$ & $\GSsame$ & $\GVopp$ & $\GSopp$ \\
\hline
{{ROWS:tbl_fm_couplings}}
\end{tabular}
\end{ruledtabular}
\end{table*}

\emph{Form C polarizes matter; form Q does not.} Table~\ref{tab:kspol} and Fig.~\ref{fig:kspol} give
the energy per particle as a function of the polarization, for the same couplings of the Lagrangian
in the two forms. In form C the energy falls with the polarization, in form Q it rises. Over
{{GEN:fm_pol_n_all}} cases (two dimensions, two sets, three strengths, four densities) the state of
lowest energy is polarized in {{GEN:fm_pol_n_cases}} cases in form C and in {{GEN:fm_pol_n_q}} cases
in form Q. The reason is Eq.~(\ref{eq:Esame}): in form C a particle is attracted by the particles of
its own half ($\GSsame<0$) and repelled by those of the other half ($\GSopp>0$), so that the
particles gather in one half and break the symmetry between the halves spontaneously; in form Q a
particle is attracted by both. (Of the {{GEN:fm_pol_n_states}} states examined,
{{GEN:fm_pol_n_missing}} states of form Q with unequal halves do not exist as solutions with
non-negative masses, at three and six times the threshold and high density; the states with equal
halves and with one filled half exist in all cases, and in form C every state exists.)

\begin{table*}
\caption{\label{tab:kspol}Energy per particle minus $m$ of matter of the total density of the Fermi
momentum $k_{F,\mathrm{tot}}=0.6\,m$ (the Fermi momentum that the density would have with equal halves),
at three polarizations, in the two forms of the $4+4$ model (units $m$). Strength: in units of the
strength at which form Q first binds matter.}
\begin{ruledtabular}
\begin{tabular}{ccccccccc}
$D$ & set & strength & C, $\mathrm{pol}_W=0$ & C, $0.5$ & C, $1$ & Q, $\mathrm{pol}_W=0$ & Q, $0.5$ & Q, $1$ \\
\hline
{{ROWS:tbl_fm_polarization}}
\end{tabular}
\end{ruledtabular}
\end{table*}

\begin{figure*}
\includegraphics[width=0.86\textwidth]{fig_forms_polarization}
\caption{\label{fig:kspol}Energy per particle of matter of a given total density against the
polarization of the two halves of the field, relative to its value for equal halves ($D=4$, set A).
Upper row: form C; lower row: form Q; left: 1.5 times, right: 3 times the strength at which form Q
first binds matter. The four curves are four total densities; the circle marks the lowest point of
each curve.}
\end{figure*}

\emph{A gas becomes self-bound matter at a definite moment}: when the strength reaches the threshold
of Table~\ref{tab:ksthr}. In form C the bound state has one half filled, in form Q equal halves, in
all eight cases; each threshold agrees with that of matter of one kind (four states with
$\GSsame,\GVsame$, or eight states with $\GStot,\GVtot$) computed by the routines of
Sec.~\ref{sec:ksmethod}, to {{GEN:fm_thr_onekind_dev}}. Matter of form C needs
{{GEN:fm_thr_ratio_d4_A}} times (set A) and {{GEN:fm_thr_ratio_d4_B}} times (set B) the strength of
form Q in $D=4$, and {{GEN:fm_thr_ratio_d3_A}} and {{GEN:fm_thr_ratio_d3_B}} times in $D=3$: on the
clock it binds later by $\Delta\td={{GENM:fm_thr_dtau_d4_A}}$ and ${{GENM:fm_thr_dtau_d4_B}}$ ($D=4$),
${{GENM:fm_thr_dtau_d3_A}}$ and ${{GENM:fm_thr_dtau_d3_B}}$ ($D=3$). Two things make the difference.
Only four of the eight internal states take part. And the attraction inside one half,
$\GSsame=\tfrac78g_0$ for set A, is weaker than the attraction $\GStot=\tfrac{15}{16}g_0$ that matter
of form Q with equal halves feels.

\begin{table*}
\caption{\label{tab:ksthr}The smallest strength $g$ of the couplings of the Lagrangian at which
uniform matter is self-bound, in the two forms of the $4+4$ model (units $m=1$), with the state just
above the threshold: its kind, the Fermi momentum $k_{F,\mathrm{tot}}$ that its density would have with
equal halves, and the masses of the particles of the two halves (for an empty half the mass that a
particle of that half would have). Last column: the largest cutoff for which the vacuum is stable at
that strength (Sec.~\ref{sec:sea}).}
\begin{ruledtabular}
\begin{tabular}{ccccccccc}
$D$ & set & form & $g\,m^{D-1}$ & state & $k_{F,\mathrm{tot}}/m$ & $m^{*}_+/m$ & $m^{*}_-/m$ & $\Lambda_{\max}/m$ \\
\hline
{{ROWS:tbl_fm_thresholds}}
\end{tabular}
\end{ruledtabular}
\end{table*}

Figure~\ref{fig:ksbind} and Table~\ref{tab:kssat} follow the saturated matter from the moment of
binding on. Its energy per particle falls monotonically and the mass of its particles falls; the
kind of the state does not change (one half filled in form C, equal halves in form Q) at any of the
101 clock readings examined up to five units of the clock after the threshold. Counted from its own
threshold, matter of form C is somewhat less strongly bound than matter of form Q at the same clock
reading. Along the table the largest eigenvalue of $\mathsf R_{\mathrm{gap}}$ is
{{GEN:fm_response_max}}. Later the saturated state of form Q reaches the eigenvalue 1, at which its
two masses can separate ({{GEN:fm_stab_list}}); that of form C does not.

\begin{table*}
\caption{\label{tab:kssat}Self-bound matter in the two forms along the deflation clock, counted for
each case from its own binding threshold: energy per particle of the saturated state minus $m$ at
four clock readings, and the mass of the particles of the filled half or halves and the Fermi
momentum $k_{F,\mathrm{tot}}$ at the clock readings 1 and 3 (units $m$).}
\begin{ruledtabular}
\begin{tabular}{ccccccccccc}
$D$ & set & form & $\td=1$ & $2$ & $3$ & $4$ & $m^{*}$ at 1 & $m^{*}$ at 3 & $k_{F,\mathrm{tot}}$ at 1 & $k_{F,\mathrm{tot}}$ at 3 \\
\hline
{{ROWS:tbl_fm_saturation}}
\end{tabular}
\end{ruledtabular}
\end{table*}

\begin{figure*}
\includegraphics[width=0.86\textwidth]{fig_forms_saturation}
\caption{\label{fig:ksbind}Energy per particle minus the mass of self-bound matter along the deflation
clock, for the same couplings of the Lagrangian in both forms. The clock is counted here from the
moment at which form Q first binds matter of the same set, so that the curves of form C begin later,
where form C first binds. Solid lines: set A; dashed lines: set B; blue: form C (one half filled);
orange: form Q (equal halves). Left: $D=4$; right: $D=3$.}
\end{figure*}

Matter of a \emph{fixed} density is carried through the same sequence (Fig.~\ref{fig:ksgas},
Table~\ref{tab:ksgas}). In form C the two halves become unequal at a definite clock reading before
the binding threshold, and one half is empty at most {{GEN:fm_gas_width}} units of the clock later:
the gas polarizes. Its pressure along the flat coordinates becomes negative after the polarization,
or, in {{GEN:fm_gas_n_no_tension}} of the 8 cases (the higher density), not at all, because the
repulsion $\GVsame$ between the particles of the filled half grows with the couplings while the
attraction saturates as the mass falls. In form Q the halves stay equal and the pressure becomes
negative in all cases: a uniform gas is then unstable against breaking up into drops of saturated
matter. Along the deflating coordinates the pressure is the interaction pressure $\pint$ alone, which
is negative as long as the scalar attraction dominates: \emph{the matter pulls on the deflating
coordinates.} Section~\ref{sec:gravity} uses this.

\begin{table}
\caption{\label{tab:ksgas}A gas of fixed total density in the two forms: the clock readings, counted
from the binding threshold of each case, at which the two halves become unequal, at which one half
is empty, and at which the pressure along the flat coordinates becomes negative. A dash: the event
does not happen for $-3\le\td\le5$.}
\begin{ruledtabular}
\begin{tabular}{ccccccc}
$D$ & set & form & $k_{F,\mathrm{tot}}/m$ & unequal & one half & $p<0$ \\
\hline
{{ROWS:tbl_fm_gas}}
\end{tabular}
\end{ruledtabular}
\end{table}

\begin{figure*}
\includegraphics[width=0.86\textwidth]{fig_forms_gas}
\caption{\label{fig:ksgas}Matter of a fixed total density while the couplings grow ($D=4$, set A; the
clock of each form is counted from its own binding threshold). (a) The polarization: in form C the
two halves become unequal at a definite clock reading and one half empties soon after; in form Q
they stay equal. (b) The pressure along the coordinates of the motion, divided by the energy
density. Solid lines: $k_{F,\mathrm{tot}}=0.3\,m$; dashed lines: $0.6\,m$.}
\end{figure*}

\subsection{Form Q for other ratios of the couplings}
\label{sec:ksformq}

Matter of form Q with equal halves is matter of one kind of particle with eight internal states,
whose two couplings are the Hartree-Fock couplings, $\GS=\GS^{\mathrm{HF}}$ and
$\GV=\GV^{\mathrm{HF}}$; it is also the matter of the $7+1$ spacetime whose fields do not depend on
the deflating coordinates. For it the equations of Sec.~\ref{sec:ksmethod} hold with
$\GS(x4)=\GS^{\mathrm{ref}}e^{\td}$ and $\GV(x4)=\GV^{\mathrm{ref}}e^{\td}$, and
Table~\ref{tab:ksthrq} gives the smallest scalar attraction that binds it for three ratios of the two
couplings (the ratio 0.2025 is that of set B). The thresholds were computed twice, by the program
and by the independent script, and agree to $10^{-9}$.

\begin{table}
\caption{\label{tab:ksthrq}Form Q with equal halves (matter of one kind with eight internal states):
the smallest scalar attraction at which uniform matter is self-bound, for $D=4$ and $D=3$ and three
ratios of the vector repulsion to the scalar attraction, with the Fermi momentum of the bound state
at the threshold (units $m=1$). Last column: the critical coupling of the vacuum for a cutoff equal
to $m$ (Sec.~\ref{sec:sea}).}
\begin{ruledtabular}
\begin{tabular}{ccccc}
$D$ & $\GV/|\GS|$ & $|\GS|\,m^{D-1}$ & $\kF/m$ & $|G_c|\,m^{D-1}$ \\
\hline
{{ROWS:tbl_ks_thresholds}}
\end{tabular}
\end{ruledtabular}
\end{table}
%%</1644>
%%<1671>
Consider now the sector of fields that do not depend on $x5,x6,x7$. Its particles have momentum
along $x0,\dots,x3$, so that $D=4$, or, if the fields do not depend on the hidden coordinate $x0$
either, along ordinary space, $D=3$.
By Eq.~(\ref{eq:couplings}) the equations of Sec.~\ref{sec:ksmethod} hold for the weighted field at
each moment with the couplings of that moment,
\begin{equation}
\GS(x4)=\GS^{\mathrm{ref}}\,e^{\td},\qquad\GV(x4)=\GV^{\mathrm{ref}}\,e^{\td},
\end{equation}
while the mass $m$ and the density of particles per unit of the volume of $x0,\dots,x3$ do not
change. The ratio $\GV/|\GS|$ is fixed. The deflation therefore carries the matter through a
one-parameter sequence of interactions of growing strength, and \emph{a gas becomes self-bound
matter at a definite moment}: when the scalar attraction reaches the threshold of
Table~\ref{tab:ksthr}.

\begin{table}
\caption{\label{tab:ksthr}The smallest scalar attraction at which uniform matter is self-bound, for
$D=4$ and $D=3$ space dimensions with momentum and three ratios of the vector repulsion to the
scalar attraction, with the Fermi momentum of the bound state at the threshold (units $m=1$). The
ratio 0.2025 is that of the Hartree-Fock couplings for $g_1=-0.1\,g_0$. Last column: the critical
coupling of the vacuum for a cutoff equal to $m$ (Sec.~\ref{sec:cjtcontact}).}
\begin{ruledtabular}
\begin{tabular}{ccccc}
$D$ & $\GV/|\GS|$ & $|\GS|\,m^{D-1}$ & $\kF/m$ & $|G_c|\,m^{D-1}$ \\
\hline
{{ROWS:tbl_ks_thresholds}}
\end{tabular}
\end{ruledtabular}
\end{table}

If the coupling at a reference time is $\GS^{\mathrm{ref}}$, the binding begins at the clock reading
$\td=\ln(|\GS|_{\mathrm{threshold}}/|\GS^{\mathrm{ref}}|)$. Figure~\ref{fig:ksbind} follows the saturated
matter from that moment on: its energy per particle falls and the mass of its particles falls, more
slowly if a vector repulsion is present. The thresholds were computed twice, by the program and by
the independent script, and agree to $10^{-9}$.

\begin{figure*}
\includegraphics[width=0.86\textwidth]{fig_ks_binding}
\caption{\label{fig:ksbind}Self-bound matter in the sector of fields that do not depend on
$x5,x6,x7$, along the deflation clock, counted from the moment at which matter first becomes
self-bound. Upper row: energy per particle of the saturated state minus $m$. Lower row: mass of its
particles. Left: $D=4$; right: $D=3$. Three ratios of the vector to the scalar coupling.}
\end{figure*}

Matter of a \emph{fixed} density is carried through the same sequence
(Fig.~\ref{fig:ksgas}). Its pressure along the flat coordinates becomes negative at a moment that
depends on the density (Table~\ref{tab:ksgas}): a uniform gas is then unstable against breaking up
into drops of saturated matter. Along the deflating coordinates the pressure is the interaction
pressure $\pint$ of Eq.~(\ref{eq:pKS}) alone, which is negative as long as the scalar attraction
dominates: \emph{the matter pulls on the deflating coordinates.} Section~\ref{sec:gravity} uses this.

\begin{table}
\caption{\label{tab:ksgas}A gas of fixed Fermi momentum in the sector of fields that do not depend on
$x5,x6,x7$: the coupling at which its pressure along the flat coordinates vanishes, and the clock
reading of that moment counted from the binding threshold of Table~\ref{tab:ksthr} (negative: before
the threshold is reached, at this higher density).}
\begin{ruledtabular}
\begin{tabular}{ccccc}
$D$ & $\GV/|\GS|$ & $\kF/m$ & $|\GS|\,m^{D-1}$ & $\td$ \\
\hline
{{ROWS:tbl_ks_gas}}
\end{tabular}
\end{ruledtabular}
\end{table}

\begin{figure*}
\includegraphics[width=\textwidth]{fig_ks_gas}
\caption{\label{fig:ksgas}A gas of fixed density in the sector of fields that do not depend on
$x5,x6,x7$ ($D=4$, $\GV=0.1\,|\GS|$), along the deflation clock counted from the binding threshold.
(a) Mass of the particles in matter. (b) Pressure along $x0,\dots,x3$ divided by the energy density.
(c) Pressure along $x5,x6,x7$, the interaction pressure alone, divided by the energy density.}
\end{figure*}
%%</1671>
